% ----------
% Tech report template
% 
% ----------
\documentclass[9pt,a4paper]{article}
\usepackage{geometry}
\usepackage{mlmodern}
\usepackage{amsmath,amsfonts,amssymb}
%\usepackage[utf8x]{inputenc}
\usepackage[T1]{fontenc}
\usepackage[english]{babel}
\usepackage[runin]{abstract}
\usepackage{titling}
\usepackage{listings}
\usepackage{booktabs}
%\usepackage{piton} % Without this enable, it might not work for the listing package. Switch to LuaLateX if you want to use this option, however. 
\usepackage{fancyhdr}
%\usepackage{helvet}
\usepackage{csquotes}
\usepackage{graphicx}
\usepackage{blindtext}
\usepackage{parskip}
\usepackage{etoolbox}
%\usepackage[scaled=0.90]{helvet} % Font 1
%\usepackage[scaled=0.85]{beramono} % Font 2
\usepackage{xcolor}
\usepackage{enumitem}
% Bibliography
\usepackage[backend=biber,style=alphabetic]{biblatex}
\addbibresource{bibb.bib} %Imports bibliography file

\input{preamble.tex}

% ----------
% Variables
% ----------

\title{RHINELAB-RINALAB - Format Standard, v3} % Full title of your tech report
\runningtitle{Organization Document} % Short title
\author{Bui Gia Khanh} % Full list of authors
\runningauthor{Amane et al.} % Short list of authors
\affiliation{} % Affiliation e.g. University or Company
\department{Core Management Division} % Department or Office
\memoid{DOC-OR-02-v2} % ID of the tech report
\theyear{2026} % year of the tech report
\mydate{May, 2026} %the date
\usepackage{tcolorbox}
\usepackage{fancyvrb}

\tcbuselibrary{minted,breakable,xparse,skins}

\definecolor{bg}{gray}{0.95}
%\DeclareTCBListing{mintedbox}{O{}m!O{}}{%
%  breakable=true,
%  listing engine=minted,
%  listing only,
%  minted language=#2,
%  minted style=default,
%  minted options={%
%    linenos,
%    gobble=0,
%    breaklines=true,
%    breakafter=,,
%    fontsize=\small,
%    numbersep=8pt,
%    #1},
%  boxsep=0pt,
%  left skip=0pt,
%  right skip=0pt,
%  left=25pt,
%  right=0pt,
%  top=3pt,
%  bottom=3pt,
%  arc=5pt,
%  leftrule=0pt,
%  rightrule=0pt,
%  bottomrule=2pt,
%  toprule=2pt,
%  colback=bg,
%  colframe=orange!70,
%  enhanced,
%  overlay={%
%    \begin{tcbclipinterior}
%    \fill[orange!20!white] (frame.south west) rectangle ([xshift=20pt]frame.north west);
%    \end{tcbclipinterior}},
%  #3}

% Python preset environment. 
\newcommand{\pyobject}[1]{\fbox{\color{red}{\texttt{#1}}}}

\NewDocumentCommand{\codeword}{v}{%
\texttt{\textcolor{blue}{#1}}%
}
% Inline python highlighting (Might not look great, but eh)
\lstset{language=Python,keywordstyle={\bfseries \color{blue}}}


% ----------
% actual document
% ----------
\begin{document}
\maketitle

\begin{abstract}
    \noindent
    Preliminary planning for the research group, denotation standard series. 

    \keywords{Standard, directive.}
\end{abstract}

\vspace{2.5cm}



\thispagestyle{firstpage}

\pagebreak

% ----------
% End of first page
% ----------

\newgeometry{} % Redefine geometries (normal margins)

\section{Specification}
\label{sec:spec}

Provide specification details, metaprogramming. None here. 

\subsection{Document code, priority}

%Provide document codes, additional document code, purpose, organization, authorship and anyone contributing to the report, identification of purpose (proposal, or research, indexing 01). Priority code, additional conditions. Does it override anything, or add to anything?

Document code administrative DOC-OR-01, version 3, Format Standard for documents, authored to serve as the reference code for all subsequent documents in line, Amane Fujimiya. Priority R1 (highest), effective immediately after release. 

\subsection{Body of issue}

Core Management Division.

\subsection{Date and Version}

Second version (v2, April 2024). Previous version has no versioning. 

\textbf{Current}: 5/11/2026 - Third version (v3). Adding tags and clarification. 

\section{Introduction}

This document outlines details on conventions of document codes. Per starter and information specified at-hand, we can say that the general format of the codes would be declared on-arrival or planning as
\begin{quote}
    [General Code] + [Specific Code] + [Sub-tags] + [Indexing 1] + [Indexing 2] + \dots
\end{quote}
and such will be clarified for all official, formal and sub-strata of actualizable documents in RHINELAB. 

\section{General code}

Generally, at this current version, artefacts generated by the laboratory will take on the following designation:
\begin{enumerate}
    \item \textbf{DOC}: Documents, general basket. 
    \item \textbf{RT}: \textbf{research topic} documents. 
    \item \textbf{OPS}: Operational objects, general basket for anything non-document and of usage. 
    \item \textbf{REP}: Report, the specific category for reports and technical notes. 
    \item \textbf{REV}: Review, specific category for comprehensive review on certain subject of interests. This encompasses literature review. 
    \item \textbf{DAT}: Data. This encompasses dataset, analytic data, test, records, structured or unstructured data form, database packets, aggregated sets, etc. 
    \item \textbf{RES}: specific category for \textbf{full research document} or artefacts of the same kind.
    \item \textbf{PED} or \textbf{FND}: Pedagogy or foundation code for document of said types. 
    \item \textbf{MAN} - large manuscript, encompassing many \textbf{RES} and \textbf{PED} document, for example. 
    \item \textbf{MOD} - categorization for modules, for example, specially modulated Python modules, work modules, analysis modules, specialized operational tools. 
    \item \textbf{MDL} - model categorization for models and other coded constructs. 
    \item \textbf{SIM} - categorization for simulation. 
    \item \textbf{TST} - categorization for stress-testing modules. 
    \item \textbf{MSC} - Miscellaneous. Filler.
\end{enumerate}

By the ordering, \textbf{DOC} is the encompassing all document, and \textbf{OPS} for codes and so on. There, \textbf{MSC} exists only for those that do not fit any basket, awaiting review and reorganization. Each document should have at most \textbf{one} (1) tag. If there exists adjacent tags, you will need to provide them in the specification section. That is, the ordering in line should be:

\begin{enumerate}
    \item \textbf{Top-tier basket}: Including codes as \textbf{MSC}, \textbf{DOC}. 
    \item \textbf{Side-documentation}: Including codes used for specific purpose, like specialized books of manuscript-size, thus \textbf{MAN}. 
    \item \textbf{Specialization}: This includes \textbf{RT,REP,REV,RES,PED,FND}. 
    \item \textbf{Action-specific or non-document}: this includes \textbf{DTA, MOD, MDL, SIM, TST}. 
\end{enumerate}

Interestingly, there have been questions about what \textbf{RT} can be used for. In short, this tag is used when the topic at hand is investigated, diagnosed of new ideas and so on, \textbf{without requiring it} to be a \textit{proposal class}. 

\subsection{Additional codes (v3)}

Version 3 introduces the following codes:

\begin{enumerate}
    \item \textbf{STM}: Statement, using on \textbf{administrative statement} or \textit{organization} purpose, alongside later administrative document format. 
    \item \textbf{DIA}: Dialogue. This is used explicitly for \textbf{long-form, condensed} and \textit{quasi-formal discussions or discourse thread}. By societal and community function, \textit{no formalization is allowed}. 
\end{enumerate}

\section{Specific tags}

Specific tags are also called \textbf{purpose tags}, which encodes utilization purpose subsequently of the main tag, for the document. Specifically, this document is categorized as \textbf{OR}, or organization since it is in consideration of the working mechanism and standards of the laboratory. Then, the few working categorization right now are,
\begin{enumerate}
    \item \textbf{OR} - organization. 
    \item \textbf{AU} - audit. 
    \item \textbf{RP} - research report. 
    \item \textbf{TE} - test (any kind and form). 
    \item \textbf{RQ} - request. 
    \item \textbf{(R)PP} - research proposals, or project proposals. Nominally use RP. 
    \item \textbf{WT} - writing (static, non-active). 
    \item \textbf{PROP} - proposals and other form of planning documents. 
\end{enumerate}
and so on. Most of the time those would be enough, and would be updated from this manuscript forward in later versions. If there exists any other tags, specify directly of those non-standard convention. There can be at most \textbf{three} (3) tags simultaneously. If there exists more than one, dash them out, i.e. \textbf{OR/AU} and so on. 

\subsection{Additional code, v3}

Specific codes are given the following addition in version 3. 

\begin{enumerate}
    \item \textbf{SM}: statement specific tag, in case for the document is with secondary interest in such. 
    \item \textbf{REA}: record analysis - for document with intention of examining or analysing records. 
    \item \textbf{RC}: records. Used to record text-based form of any given interest, internal to the lab \textbf{only} and in public space \textit{only}, with respect to the laboratory itself. 
\end{enumerate}

\subsection{Security code}

Security code is simply there, as of its purpose for now and the projection in between, to be the \textit{access code} that would be required per sensitive or administrative information, which is not to be disclosed, within affects. 

\begin{enumerate}
    \item \textbf{S3}: Absolute security, equal to private, and administrative bodies-based access points, i.e. Core Management Division. 
    \item \textbf{S2}: Private internal level, equal to lab private access, members of the laboratory, only to director can access such. 
    \item \textbf{S1}: Private internal level, equal to lab private access, all member of the laboratory, in official full member can access such. 
    \item \textbf{S0}: public, full exposure. 
\end{enumerate}

The determination of those security codes are to be investigated under scrutiny. Often, such information are to be considered \textit{pre-applied}, however in case of which the reverse are to happen, identification and audit-pipeline review are needed. Usually, all documents are S0 for public domain, and such are not required to list security code. 

\section{Indexing}

Indexing is simple. We have the three main indexing partitions:
\begin{enumerate}
    \item \textbf{Alphabetical} (singular, uppercase) - from A to Z, 26 indexes. 
    \item \textbf{Alphabetical} (singular, lowercase) - from a to z, 26 indexes.
    \item \textbf{Roman} (singular, uppercase) - from I to C. 100 indexes. 
    \item \textbf{Roman} (singular, lowercase) - from i to li. 51 indexes. 
    \item \textbf{Base-10 numerics} - from 0 to arbitrary $n$. Proliferation allowed. 
\end{enumerate}

They are indexes that you use for various purposes, but also a type of \textbf{range enforcement} and anti-proliferation. Choose under constraints of the \textit{scope} of how much you want the content to spread out. 


\section{Specification template}

In R1D1 document (DOC-OR-01), the specification section is listed similar to the following picture,
\begin{figure}[htb]
    \centering
    \includegraphics[width=0.6\textwidth]{img/R1_example.png}
    \caption{Example specification of document.}
\end{figure}

Using this, we require at least three fields:
\begin{enumerate}
    \item Document code, priority. 
    \item Body of issues. 
    \item Date and version. 
\end{enumerate}

Priority answers if this is urgent or not, base on \texttt{R}-code from zero (0), one (1) to five (5), with zero is the highest. Body of issue follow from the document. 


% ----------
% Bibliography
% ----------

% \bibliography{../biblio.bib}
% \bibliographystyle{abbrv}

%\appendix



%\blindenumerate


%\medskip

%\printbibliography %Prints bibliography


\end{document}

% ----------
